% !TEX program = xelatex
% !TeX encoding = UTF-8
\documentclass[12pt,a4paper]{article}
\usepackage{fontspec}
\usepackage{polyglossia}
\setmainlanguage{arabic}
\setotherlanguage{english}

% خطوط عربية
\newfontfamily\arabicfont{Amiri}[Script=Arabic, Language=Arabic]
\newfontfamily\arabicfontsf{Amiri}[Script=Arabic, Language=Arabic]
\newfontfamily\arabicfonttt{Amiri}[Script=Arabic, Language=Arabic]

% خطوط لاتينية
\setmainfont{Times New Roman}[Ligatures=TeX]
\setsansfont{Arial}
\setmonofont{Consolas}[Scale=MatchLowercase]

% حزم الرياضيات والتنسيق
\usepackage{amsmath,amssymb,amsthm,mathtools}
\usepackage{geometry}
\geometry{left=2.5cm,right=2.5cm,top=2.5cm,bottom=2.5cm}
\usepackage{booktabs}
\usepackage{hyperref}
\usepackage{url}
\usepackage{enumitem}
\usepackage{float}
\usepackage{xcolor}
\usepackage{titlesec}
\usepackage{fancyhdr}

% تعريف الألوان
\definecolor{gold}{RGB}{255,215,0}
\definecolor{darkblue}{RGB}{0,0,139}
\definecolor{skyblue}{RGB}{0,102,204}

% مكروبات YUCT
\newcommand{\Keff}{K_{\mathrm{eff}}}
\newcommand{\kappac}{\kappa_c}
\newcommand{\eps}{\varepsilon}
\newcommand{\df}{d_{\!f}}
\newcommand{\dint}{\ensuremath{D\!+\!I\!\cdot\!R}}
\newcommand{\Sodd}{S_{\mathrm{odd}}}
\newcommand{\Seven}{S_{\mathrm{even}}}

% بيئات النظريات – عربية
\newtheorem{theorem}{نظرية}[section]
\newtheorem{definition}[theorem]{تعريف}
\newtheorem{proposition}[theorem]{اقتراح}
\newtheorem{conjecture}[theorem]{تخمين}
\theoremstyle{remark}
\newtheorem{remark}[theorem]{ملاحظة}

\hypersetup{
	colorlinks=true,
	linkcolor=blue,
	citecolor=blue,
	urlcolor=blue,
}

% تنسيق العناوين (تم حذف السطر المسبب للخطأ)
\titleformat{\section}{\LARGE\bfseries\color{darkblue}}{\arabic{section}}{1em}{}
\titleformat{\subsection}{\Large\bfseries\color{darkblue}}{\arabic{subsection}}{1em}{}

% تذييل الصفحة
\fancypagestyle{firstpage}{%
	\fancyhf{}%
	\renewcommand{\headrulewidth}{0pt}%
	\renewcommand{\footrulewidth}{0pt}%
	\fancyfoot[C]{%
		\begin{minipage}[t]{\textwidth}
			\centering
			\textcolor{gold}{\rule{\textwidth}{1.6pt}}\\[5pt]
			{\small \textsuperscript{1}أبحاث ياكوشيف، \url{https://yuct.org/}} \hfill
			{\small بروتوكول YPSDC: \url{https://ypsdc.com/}}\\[-2pt]
			\textcolor{gold}{\rule{\textwidth}{1.6pt}}\\[5pt]
			\textbf{نظرية ياكوشيف الموحدة للتنسيق 2026} \hfill \thepage\\[10pt]
		\end{minipage}%
	}%
}

\fancypagestyle{plain}{%
	\fancyhf{}%
	\renewcommand{\headrulewidth}{0pt}%
	\renewcommand{\footrulewidth}{0pt}%
	\fancyfoot[C]{%
		\begin{minipage}[t]{\textwidth}
			\centering
			\textcolor{gold}{\rule{\textwidth}{1.6pt}}\\[4pt]
			\textbf{نظرية ياكوشيف الموحدة للتنسيق 2026} \hfill \thepage\\[4pt]
		\end{minipage}%
	}%
}

\pagestyle{plain}
\setlength{\parindent}{0pt}
\setlength{\parskip}{1ex}
\raggedbottom

\makeatletter
\let\@ensure@LTR\relax
\makeatother

% ============================================================
% بداية المستند
% ============================================================
\begin{document}
	
	\begin{titlepage}
		\centering
		\vspace*{1cm}
		
		\Huge\textbf{الملحق F (المنقح): الديناميكيات الكسيرية للقواميس الاقتصادية ونظرية التنسيق للنقود}
		
		\vspace{0.5cm}
		\Large\textbf{إطار نظري لأسعار الصرف والمناطق النقدية والحدود الاقتصادية ذاتية التنظيم}
		
		\vspace{1.5cm}
		
		\Large\textbf{أليكسي ف. ياكوشيف\textsuperscript{1}}
		
		نظرية YUCT: \url{https://yuct.org/}\\
		بروتوكول YPSDC: \url{https://ypsdc.com/}
		
		\vspace{2cm}
		\textbf{DOI:} \url{https://doi.org/10.5281/zenodo.18444598}
		
		\vspace{1cm}
		
		\large
		يونيو 2026 \\ \large الإصدار 3.0 (يحل محل الملحق F السابق)
		
		\vspace{2cm}
		
		\begin{minipage}{0.9\textwidth}
			\begin{abstract}
				\noindent
				يقدم هذا الملحق مراجعة جذرية لمنهج YUCT في الاقتصاد. نتخلى عن جميع المحاولات السابقة لتركيب الناتج المحلي الإجمالي أو تقلبات السوق بشكل تجريبي مع الأس العالمي \(\beta = 0.67\)، وبدلاً من ذلك نطور إطاراً نظرياً من المبادئ الأولى يقوم على أنطولوجيا التنسيق. يُنمذج الاقتصاد كشبكة من \textbf{قواميس اقتصادية} متفاعلة (شركات، قطاعات، دول)، يتميز كل منها بكفاءة تنسيق \(K_{\mathrm{eff}}\). نبين أن الحدود بين أي قاموسين تبدي ديناميكيات كسيرية مشابهة لمجموعة ماندلبروت: متشابهة ذاتياً، معقدة، ومحكومة بقانون الخطأ العام \(\varepsilon = \kappa_c \alpha K_{\mathrm{eff}}^{-\beta}\). يُعرَّف المال بأنه \textbf{معامل التنسيق} الذي ينظم التنشيط عبر الحدود القاموسية؛ وبالتالي تتحدد أسعار الصرف بنسبة كفاءات التنسيق. ينتج الإطار عدة تنبؤات كمية قابلة للتفنيد: (1) يجب أن يتجمع أس هيرست لسلاسل أسعار الصرف عالية التردد حول \(H \approx 1/\beta \approx 1.5\) في المقاييس الزمنية الطويلة؛ (2) يجب أن يظهر الطيف متعدد الكسير للعوائد المالية تبايناً يسارياً ذا شكل عالمي تحدده الحلقة الجبرية الأولية؛ (3) يتناسب التقلب \(\sigma\) لسعر الصرف مع فجوة التنسيق \(\Delta K_{\mathrm{eff}}\) بين القاموسين كـ \(\sigma \propto (\Delta K_{\mathrm{eff}})^{-0.67}\). لا يتم إجراء أي تركيب تجريبي في هذا الملحق؛ تُصاغ جميع التنبؤات بشكل يمكن اختباره بمجرد توفر وكلاء موثوقين لـ \(K_{\mathrm{eff}}\) في الأنظمة الاقتصادية. يحوّل هذا العمل اقتصاد YUCT من مجموعة ارتباطات مشكوك فيها إلى برنامج بحثي صارم وقابل للتفنيد.
			\end{abstract}
		\end{minipage}
		
		\vspace{1cm}
		
		\noindent
		\small\textbf{الكلمات المفتاحية:} YUCT، كفاءة التنسيق، القواميس الاقتصادية، أسعار الصرف، الحدود الكسيرية، مجموعة ماندلبروت، أس هيرست، التحليل متعدد الكسير، المال كمعامل تنسيق.
		
		\vspace{0.5cm}
		
		\noindent
		\textcopyright\ 2026 أبحاث ياكوشيف. جميع الحقوق محفوظة.
	\end{titlepage}
	
	\tableofcontents
	\newpage
	
	% ============================================================
	% SECTION 1: INTRODUCTION AND MOTIVATION
	% ============================================================
	\section{مقدمة: لماذا بداية جديدة ضرورية}
	
	حاولت الإصدارات السابقة من هذا الملحق إثبات أن السلاسل الزمنية الاقتصادية الكلية — تقلبات نمو الناتج المحلي الإجمالي، وأخطاء التنبؤ بالعملات المشفرة، ومؤشرات المخاطر الجيوسياسية — تتبع قانون القياس العام \(\varepsilon \propto K_{\mathrm{eff}}^{-0.67}\) لـ YUCT. كشف التحقق المستقل اللاحق أن العديد من جداول البيانات والانحدارات المذكورة كانت إما خاطئة أو اعتمدت على اختيار لاحق للحالات الملائمة. يعترف هذا المنقح بالكامل بهذه القصور ويسحب جميع الادعاءات التجريبية السابقة.
	
	لم يكن الخطأ الأساسي في إطار YUCT نفسه، بل في محاولة فرض تركيب تجريبي مباشر دون تطوير نموذج نظري مناسب أولاً لما هو "القاموس الاقتصادي" فعلاً وكيف يمكن قياس كفاءة تنسيقه. الاقتصاد ليس فيزياء الجسيمات: الوكلاء الاقتصاديون ليسوا جسيمات نقطية، وكتلهم ليست ثوابت أساسية. لذلك، كان البحث عن سلم كتلي بسيط في البيانات الاقتصادية مضللاً منذ البداية.
	
	يقدم هذا الملحق نهجاً مختلفاً جذرياً. بدلاً من معاملة الاقتصاد كمجموعة من الأشياء التي يجب أن تقع أحجامها على طيف منفصل، نعامله كـ \textbf{شبكة ديناميكية من القواميس المتفاعلة}، التي تمثل حدودها الموضع الأساسي للظواهر المثيرة للاهتمام. أسعار الصرف، وأسعار الفائدة، والتقلبات المالية ليست خصائص لوكلاء معزولين، بل \textbf{خصائص ناشئة للحدود} بين القواميس. وهندسة هذه الحدود، كما نزعم، هي كسيرية ومتشابهة ذاتياً، وتحكمها نفس الأس العالمي \(\beta = 0.67\) الذي يظهر في جميع مجالات YUCT الأخرى.
	
	لا يتم إجراء أي تركيبات تجريبية في هذا الملحق. جميع العبارات الكمية مستمدة رياضياً من بديهيات YUCT وتُقدَّم كتنبؤات قابلة للتفنيد للبحث المستقبلي.
	
	% ============================================================
	% SECTION 2: ECONOMIC DICTIONARIES AND THEIR BOUNDARIES
	% ============================================================
	\section{القواميس الاقتصادية وحدودها}
	
	\subsection{تعريف القاموس الاقتصادي}
	
	\begin{definition}[القاموس الاقتصادي]
		\textbf{القاموس الاقتصادي} \(\mathcal{D}_i\) هو مجموعة مهيكلة من الإجراءات الممكنة، والعقود، وبروتوكولات الإنتاج، وقواعد التسعير، والأعراف المؤسسية التي يمكن للوكيل الاقتصادي (شركة، أسرة، قطاع، دولة) تفعيلها استجابةً لمؤشرات خارجية. يُقاس حجم القاموس بأنتروبيا شانون \(H(\mathcal{D}_i)\).
	\end{definition}
	
	أمثلة:
	\begin{itemize}
		\item يشمل قاموس الشركة كتالوج منتجاتها، وعقود التوريد، وخوارزميات التسعير، والروتين التنظيمي الداخلي.
		\item يشمل قاموس الاقتصاد الوطني نظامه القانوني، وقانون الضرائب، والاتفاقيات التجارية، وإطار السياسة النقدية، والبنية التحتية المالية.
	\end{itemize}
	
	كفاءة التنسيق للقاموس \(i\) هي
	\begin{equation}
		K_{\mathrm{eff}}^{(i)} = \frac{H(\mathcal{A}_i)}{H(\kappa_i)},
		\label{eq:keff_econ}
	\end{equation}
	حيث \(\mathcal{A}_i\) هي مجموعة الإجراءات التي يمكن للقاموس إنتاجها، و \(\kappa_i\) هو المؤشر (إشارة سعر، توجيه تنظيمي، أمر سوق) الذي يفعلها.
	
	\subsection{الحدود بين قاموسين}
	
	\begin{definition}[حدود القاموس]
		\textbf{الحدود} \(\partial\mathcal{D}_{ij}\) بين قاموسين اقتصاديين \(\mathcal{D}_i\) و \(\mathcal{D}_j\) هي مجموعة جميع المعاملات والعقود وتبادلات المعلومات التي يُفعِّل فيها مؤشر ينشأ في أحد القاموسين مدخلاً في الآخر.
	\end{definition}
	
	الحدود هي الموضع الطبيعي لـ:
	\begin{itemize}
		\item أسعار الصرف (للقواميس الوطنية)،
		\item تكاليف المعاملات وفروق الأسعار (للقواميس على مستوى الشركات)،
		\item عدم تناظر المعلومات والاختيار السلبي.
	\end{itemize}
	
	\begin{proposition}[معامل التنسيق على الحدود]
		تُقاس كفاءة التنسيق عبر الحدود \(\partial\mathcal{D}_{ij}\) بمعامل عديم الأبعاد \textbf{معامل التنسيق} \(\chi_{ij}\)، المُعرَّف كـ
		\begin{equation}
			\chi_{ij} = \frac{K_{\mathrm{eff}}^{(j)}}{K_{\mathrm{eff}}^{(i)}} \cdot \rho_{ij},
			\label{eq:chi}
		\end{equation}
		حيث \(\rho_{ij}\) هو عامل رنين يحدد مدى توافق القاموسين (معايير مشتركة، لغة مشتركة، تنسيق قانوني، إلخ).
	\end{proposition}
	
	\textbf{المال كمعامل تنسيق.} في حالة اقتصادين وطنيين بعملات مستقلة، فإن معامل التنسيق \(\chi_{ij}\) هو بالضبط \textbf{سعر الصرف} (وحدات العملة \(j\) لكل وحدة عملة \(i\))، حتى ثابت تطبيع. إن \(K_{\mathrm{eff}}^{(j)}\) الأعلى بالنسبة لـ \(K_{\mathrm{eff}}^{(i)}\) يعني أن القاموس \(j\) أكثر "قيمة" — تُفعَّل مدخلاته بشكل أكثر موثوقية — وبالتالي فإن عملته تطلب علاوة.
	
	% ============================================================
	% SECTION 3: FRACTAL DYNAMICS OF THE BOUNDARY
	% ============================================================
	\section{الديناميكيات الكسيرية للحدود}
	
	\subsection{القياس مع مجموعة ماندلبروت}
	
	تُعرَّف مجموعة ماندلبروت \(\mathcal{M}\) بالخريطة التكرارية \(z_{n+1} = z_n^2 + c\)، مع \(z_0 = 0\). النقاط \(c\) التي تظل المتتالية فيها محدودة تنتمي إلى \(\mathcal{M}\)؛ والنقاط التي تتباعد تقع خارجها. الحدود \(\partial\mathcal{M}\) هي كسيرية ذات تعقيد لا نهائي، وتظهر تشابهاً ذاتياً على جميع المقاييس.
	
	نقترح نظيراً اقتصادياً. دع قاموسين \(\mathcal{D}_i\) و \(\mathcal{D}_j\) يتفاعلان عبر سلسلة من المعاملات المفهرسة بخطوات زمنية متقطعة \(n\). تُوصف حالة التنسيق بعد \(n\) معاملة بمتغير مركب \(z_n \in \mathbb{C}\)، يمثل جزءاه الحقيقي والتخيلي درجة التماثل في بُعدين متعامدين (مثل تزامن الأسعار والوفاء التعاقدي). قاعدة التحديث هي
	\begin{equation}
		z_{n+1} = z_n^2 + \chi_{ij},
		\label{eq:mandel_econ}
	\end{equation}
	حيث \(\chi_{ij} \in \mathbb{C}\) هو معامل التنسيق (المُركَّب) من المعادلة~(\ref{eq:chi}). المقدار \(|\chi_{ij}|\) يقيس قوة الاقتران؛ أما طوره فيقيس إزاحة الطور بين دورات تفعيل القاموسين.
	
	\begin{conjecture}[التطابق الاقتصادي مع ماندلبروت]
		يمكن لزوج القواميس \((\mathcal{D}_i, \mathcal{D}_j)\) أن يحافظ على تنسيق محدود ومستقر (أي تبقى المتتالية \(\{z_n\}\) محدودة) إذا وفقط إذا كان \(\chi_{ij}\) يقع ضمن مجموعة ماندلبروت المعممة \(\mathcal{M}_{\mathrm{econ}}\). تفصل الحدود \(\partial\mathcal{M}_{\mathrm{econ}}\) مناطق التنسيق المستقر (أسعار صرف ثابتة، تدفقات تجارية يمكن توقعها) عن مناطق التباعد الفوضوي (أزمات العملات، التضخم المفرط، انهيار التجارة). هندسة هذه الحدود كسيرية، بتراكيب متشابهة ذاتياً على جميع المقاييس الزمنية.
	\end{conjecture}
	
	\subsection{القياس العالمي على الحدود}
	
	لأن الديناميكيات على الحدود تخضع لنفس قانون خطأ التنسيق الذي يسري في كل مكان في YUCT، يجب أن تتبع الخصائص الإحصائية لأسعار الصرف قرب الحدود قوانين قياس عالمية. تحديداً، ليكن \(\sigma_{ij}(\Delta t)\) هو تقلب سعر الصرف بين القاموسين \(i\) و \(j\) المقاس على أفق زمني \(\Delta t\). عندئذ، من قانون الخطأ العالمي \(\varepsilon \propto K_{\mathrm{eff}}^{-\beta}\)، نحصل على:
	
	\begin{proposition}[قياس تقلب أسعار الصرف]
		يتناسب تقلب معامل التنسيق \(\chi_{ij}\) مع فجوة التنسيق \(\Delta K_{\mathrm{eff}} = |K_{\mathrm{eff}}^{(i)} - K_{\mathrm{eff}}^{(j)}|\) كـ
		\begin{equation}
			\sigma(\chi_{ij}) \propto (\Delta K_{\mathrm{eff}})^{-\beta} = (\Delta K_{\mathrm{eff}})^{-0.67}.
			\label{eq:vol_scaling}
		\end{equation}
		تظهر أزواج القواميس ذات كفاءات تنسيق متشابهة جداً (فجوة صغيرة \(\Delta K_{\mathrm{eff}}\)) تقلباً عالياً وفوضويّاً؛ أما الأزواج ذات الفجوة الكبيرة فتبدي تقلباً منخفضاً ومستقراً.
	\end{proposition}
	
	هذا تنبؤ قابل للتفنيد. بمجرد تطوير وكلاء موثوقين لـ \(K_{\mathrm{eff}}^{(i)}\) (مثلاً من مؤشرات جودة المؤسسات، أو مقاييس تعقيد التجارة، أو الإنتروبيا المعلوماتية للقوانين)، يمكن اختبار المعادلة~(\ref{eq:vol_scaling}) على بيانات لوحية عبر البلدان.
	
	\subsection{الخصائص الكسيرية للسلاسل الزمنية المالية}
	
	من الحقائق التجريبية المعروفة في الاقتصاد الفيزيائي أن السلاسل الزمنية المالية تبدي خصائص متعددة الكسير: تجمع التقلب، ذاكرة طويلة، وتوزيعات عوائد ذات ذيول ثقيلة. ضمن YUCT، هذه ليست مصادفات بل نتائج ضرورية للديناميكيات الكسيرية على حدود القاموس.
	
	\begin{proposition}[أس هيرست على الحدود]
		بالنسبة لزوج من القواميس يعمل بالقرب من الحدود الحرجة للتنسيق المستقر، فإن أس هيرست \(H\) لسلسلة سعر الصرف يحقق
		\begin{equation}
			H \to \frac{1}{\beta} = \frac{3}{2} \quad \text{عندما} \quad \Delta t \to \infty.
			\label{eq:hurst}
		\end{equation}
		في المقاييس الزمنية الأقصر، يبدي \(H\) تقاطعاً إلى قيم أدنى، مما يعكس الطبيعة متعددة الكسير للحدود.
	\end{proposition}
	
	هذا التنبؤ قابل للاختبار بتحليل R/S القياسي أو DFA على بيانات أسعار الصرف عالية التردد. القيمة \(H = 1.5\) هي سمة لعملية ذات ذاكرة طويلة جداً، تقترب من نقطة تفرد — وهو السلوك الملاحظ بالضبط في الأسواق المالية قبل الأزمات الكبرى (مثل أزمة 2008).
	
	\subsection{التنظيم الكسيري الداخلي لقاموس واحد}
	
	بينما تكون الحدود بين القواميس كسيرية وفوضوية، فإن البنية الداخلية لقاموس واحد حسن التنسيق تبدي تنظيماً هرمياً متقطعاً. مدخلات القاموس ليست موزعة بشكل منتظم، بل تشكل شجرة متشعبة متشابهة ذاتياً من القواعد والعقود والبروتوكولات. أحجام الشركات، وتوزيع الدخول، والتسلسل الهرمي للأنظمة الحضرية كلها مظاهر لهذه البنية الكسيرية الداخلية.
	
	على عكس المحاولة السابقة الفاشلة لوضع اقتصادات بأكملها على سلم كتلي، نتنبأ الآن أنه \textbf{ضمن} قاموس واحد، يجب أن تتبع أحجام الوحدات الفرعية (الشركات، المدن، فئات الدخل) قانون قوى يرتبط أسّه بالبعد الكسيري \(d_f\) و \(\beta\) العام. إن قانون زيبف المعروف للمدن (أس \(\approx 1\)) وقانون باريتو للدخول (أس \(\approx 1.5\)) ليسا غرائب معزولة، بل إسقاطات مختلفة لنفس الهندسة التنسيقية الأساسية.
	
	% ============================================================
	% SECTION 4: MATHEMATICAL FRAMEWORK
	% ============================================================
	\section{الإطار الرياضي: لاغرانجيان التنسيق على الحدود}
	
	لجعل الصورة النوعية أعلاه دقيقة رياضياً، نضمّنها في لاغرانجيان YUCT ذي الـ 19 بُعداً. ليُوصف القطاع الاقتصادي بحقل \(\Psi_{72}(x^\mu, y^m)\)، حيث \(x^\mu\) هي الإحداثيات الزمكانية الأربعة و \(y^m\) هي الإحداثيات الداخلية (المتراصة) الـ 15. يقابل القاموس \(\mathcal{D}_i\) قيمة توقع فراغية معينة \(\langle \Psi_{72} \rangle_i\)، والحدود بين قاموسين هي حل جدار مجال يتوسط بين \(\langle \Psi_{72} \rangle_i\) و \(\langle \Psi_{72} \rangle_j\).
	
	الفعل الفعال على الحدود هو
	\begin{equation}
		S_{\mathrm{boundary}} = \int d^4x \sqrt{-g} \left[ \frac{1}{2} \kappa (\nabla \chi)^2 + V(\chi) \right],
		\label{eq:action}
	\end{equation}
	حيث \(\chi(x)\) هو معامل التنسيق المحلي (حقل سعر الصرف)، و \(\kappa\) هي صلابة الحدود، والجهد \(V(\chi)\) معطى بـ
	\begin{equation}
		V(\chi) = V_0 \left( e^{\lambda(\chi - \chi_0)} - \lambda(\chi - \chi_0) - 1 \right), \qquad \lambda = \frac{1}{\beta} = \frac{3}{2}.
		\label{eq:potential}
	\end{equation}
	هذا هو نفس الجهد الذي يحكم حقل التنسيق القياسي في القطاع الجذبوي (الملحق G)، مما يضمن العالمية.
	
	معادلة الحركة الكلاسيكية لـ \(\chi\) هي
	\begin{equation}
		\kappa \Box \chi - V'(\chi) = 0,
		\label{eq:eom}
	\end{equation}
	وتحقق تقلباته الإحصائية
	\begin{equation}
		\langle (\delta \chi)^2 \rangle \propto \int \frac{d^4k}{(2\pi)^4} \frac{1}{\kappa k^2 + V''(\chi_0)} \propto \chi_0^{-2\beta},
		\label{eq:fluctuations}
	\end{equation}
	معيداً إنتاج قانون قياس التقلب المعادلة~(\ref{eq:vol_scaling}).
	
	% ============================================================
	% SECTION 5: FALSIFIABLE PREDICTIONS
	% ============================================================
	\section{تنبؤات قابلة للتفنيد}
	
	يقدّم الإطار المطوّر أعلاه عدة تنبؤات كمية يمكن اختبارها دون أي تركيب بيانات:
	
	\begin{enumerate}[label=(\arabic*)]
		\item \textbf{أس هيرست لأسعار الصرف.} بالنسبة لأزواج العملات الرئيسية ذات التعويم الحر (مثل EUR/USD، USD/JPY)، يجب أن يقع أس هيرست طويل الأجل المحسوب من البيانات اليومية على مدى 20 عاماً في المجال \(1.3 \le H \le 1.7\). يجب رفض الفرضية الصفرية للتمويل السائد (المشي العشوائي، \(H = 0.5\)) بثقة عالية.
		
		\item \textbf{قياس التقلب–فجوة التنسيق.} بمجرد بناء وكيل موثوق لـ \(K_{\mathrm{eff}}\) (مثلاً من مؤشرات الحوكمة العالمية للبنك الدولي المدمجة مع بيانات تعقيد التجارة)، يجب أن يتناسب التقلب المقطعي لأسعار الصرف مع الفرق المطلق \(\Delta K_{\mathrm{eff}}\) كـ \(\sigma \propto (\Delta K_{\mathrm{eff}})^{-0.67}\). يمكن تقدير الأس بالانحدار الخطي في إحداثيات لوغاريتمية؛ التنبؤ هو أن الميل هو \(-0.67 \pm 0.10\).
		
		\item \textbf{تباين الطيف متعدد الكسير.} يجب أن يبدي طيف التفرد \(f(\alpha)\) لعوائد أسعار الصرف تبايناً يسارياً بمعامل تباين \(A_\alpha \approx 0.30 \pm 0.05\)، مما يعكس هيمنة الانهيارات المنسقة الكبيرة على الارتفاعات المنسقة الكبيرة. يمكن اختبار ذلك ببرامج MF-DFA القياسية على البيانات المتاحة للعموم.
		
		\item \textbf{التباطؤ الحرج قبل أزمات العملات.} عندما يقترب زوج عملات من حد حرج (أي عندما \(\Delta K_{\mathrm{eff}} \to 0\))، يجب أن يتباعد زمن الارتباط الذاتي للتقلب. يتنبأ هذا بأنه قبل أزمة عملة كبرى، يجب أن نلاحظ زيادة منتظمة في أس هيرست وتضيقاً في الطيف متعدد الكسير، مما يوفر إشارة إنذار مبكر.
	\end{enumerate}
	
	\textbf{حالة التنبؤات:} جميع التنبؤات مستمدة رياضياً من بديهيات YUCT دون أي تركيب للبيانات الاقتصادية. وهي مقدمة بشكل يمكن اختباره مباشرة من قبل باحثين مستقلين. تأكيد أي واحد منها سيوفر دليلاً قوياً لنظرية التنسيق للنقود؛ وتفنيد جميع الأربعة سيتطلب مراجعة جوهرية للإطار.
	
	% ============================================================
	% SECTION 6: OPEN PROBLEMS AND PATH FORWARD
	% ============================================================
	\section{مسائل مفتوحة والطريق إلى الأمام}
	
	\begin{enumerate}[label=(\roman*)]
		\item \textbf{بناء وكيل موثوق لـ \(K_{\mathrm{eff}}\) للدول.} المهمة الأكثر إلحاحاً هي تطوير طريقة لتقدير كفاءة التنسيق لاقتصاد وطني من البيانات المتاحة للعموم. تشمل المرشحات: إنتروبيا القانون، تعقيد سلة الصادرات (مؤشر هيدالغو–هاوسمان)، جودة المؤسسات (مؤشرات الحوكمة العالمية)، ومقاييس معلوماتية للبنية الدقيقة للسوق. يمكن عندئذ استقراء مؤشر مركب تمت معايرته على مجموعة صغيرة من الدول المرجعية إلى اللوحة الكاملة.
		
		\item \textbf{المحاكاة العددية لخريطة ماندلبروت الاقتصادية.} يجب دراسة الخريطة \(z_{n+1} = z_n^2 + \chi\) بشكل منهجي. لأي قيم \(\chi\) تظل المتتالية محدودة؟ ما هو البعد الكسيري للحدود؟ كيف يؤثر إضافة الضوضاء (درجة حرارة منتهية) على البنية؟ يمكن معالجة هذه الأسئلة بنصوص بايثون بسيطة وقد تكشف عن روابط كمية غير متوقعة مع ديناميكيات سعر الصرف الحقيقية.
		
		\item \textbf{الربط بالاقتصاد الفيزيائي الراسخ.} تتداخل التنبؤات المذكورة في القسم~5 بشكل كبير مع النتائج الموجودة في الاقتصاد الفيزيائي (تعدد الكسير للعوائد المالية، الذاكرة الطويلة في التقلب، الذيول ذات قانون القوى). سيوفر استعراض منهجي لهذه الأدبيات، معاد تفسيره من خلال عدسة YUCT، اختباراً فورياً للإطار دون الحاجة إلى جمع بيانات جديدة.
	\end{enumerate}
	
	% ============================================================
	% SECTION 7: CONCLUSION
	% ============================================================
	\section{الاستنتاج}
	
	لقد أعدنا هيكلة منهج YUCT في الاقتصاد بالكامل. يتخلى الإطار الجديد عن الاستراتيجية الفاشلة لإجبار المجاميع الاقتصادية على السلم الكتلي، ويطور بدلاً من ذلك نظرية من المبادئ الأولى للقواميس الاقتصادية وحدودها. يُعرَّف المال بأنه معامل التنسيق المنظّم للتنشيط عبر القواميس، وتُظهر أسعار الصرف أنها محكومة بنفس الأس العالمي \(\beta = 0.67\) الذي يظهر في كل مكان آخر في YUCT. التنبؤات الناتجة كمية وقابلة للتفنيد ومستقلة عن أي تركيب بيانات لاحق. يحول هذا اقتصاد YUCT من عبء إلى برنامج بحثي واعد وصارم.
	
	\section*{إخلاء المسؤولية}
	
	هذا الملحق هو عمل نظري. لا تُقدَّم أي ادعاءات تجريبية، ولا يُقصد به أي توصيات استثمارية أو سياسية. جميع التنبؤات الكمية معروضة للتدقيق العلمي وقد تُفنَّد ببيانات مستقبلية.
	
	\section*{توفر البيانات}
	
	لم تُنشأ بيانات جديدة لهذه الدراسة النظرية. جميع المصادر المذكورة متاحة للعموم.
	
	\begin{thebibliography}{99}
		\bibitem{AppendixA} أ.ف.ياكوشيف، ``الملحق A: لاغرانجيان YUCT ذو الـ 19 بُعداً,'' في \emph{YUCT}، Zenodo (2026).
		\bibitem{AppendixL} أ.ف.ياكوشيف، ``الملحق L: قياس الخطأ الكسيري للتنسيق,'' في \emph{YUCT}، Zenodo (2026).
		\bibitem{AppendixG} أ.ف.ياكوشيف، ``الملحق G: الجاذبية كتوتر هندسي,'' في \emph{YUCT}، Zenodo (2026).
		\bibitem{AppendixX} أ.ف.ياكوشيف، ``الملحق X: تعميم نظرية شانون للمعلومات,'' في \emph{YUCT}، Zenodo (2026).
		\bibitem{Mandelbrot1982} B.B. ماندلبروت، \emph{الهندسة الكسيرية للطبيعة}، W.H. فريمان (1982).
		\bibitem{Pareto1896} V. باريتو، \emph{دورة الاقتصاد السياسي}، Rouge (1896).
		\bibitem{Zipf1949} G.K. زيبف، \emph{السلوك البشري ومبدأ أقل جهد}، أديسون-ويسلي (1949).
		\bibitem{Hidalgo2009} C.A. هيدالغو، R. هاوسمان، ``لبنات التعقيد الاقتصادي،'' \emph{PNAS} \textbf{106}، 10570–10575 (2009).
		\bibitem{Kantelhardt2002} J.W. كانتلهارت وآخرون، ``تحليل التقلبات متعدد الكسير المنزوع الاتجاه للسلاسل الزمنية غير الثابتة،'' \emph{Physica A} \textbf{316}، 87–114 (2002).
	\end{thebibliography}
	
\end{document}